% Этот пакет позволит ссылаться на нумерованные объекты во внешних документах
% в нашем случае -- на счётчики страниц, таблиц, рисунков и т.д.
\usepackage{xr} %
\usepackage{refcount}

%%% Микротипографика %%%
%\ifnumequal{\value{draft}}{0}{% Только если у нас режим чистовика
%    \usepackage[final]{microtype}[2016/05/14] % улучшает представление букв и слов в строках, может помочь при наличии отдельно висящих слов
%}{}
\usepackage{fr-longtable}    %ради \endlasthead
\usepackage{tablefootnote}

% Листинги с исходным кодом программ
\usepackage{fancyvrb}
\usepackage{listings}
\lccode`\~=0\relax %Без этого хака из-за особенностей пакета listings перестают работать конструкции с \MakeLowercase и т. п. в (xe|lua)latex

% Русская традиция начертания греческих букв
\usepackage{upgreek} % прямые греческие ради русской традиции



\ifnumequal{\value{draft}}{1}{% Черновик
	\IfFileExists{.git/gitHeadInfo.gin}{
		\usepackage[mark,pcount]{gitinfo2}
		\renewcommand{\gitMark}{rev.\gitAbbrevHash\quad\gitCommitterEmail\quad\gitAuthorIsoDate}
		\renewcommand{\gitMarkFormat}{\rmfamily\color{Gray}\small\bfseries}
	}{}
}{}

\usepackage{nicefrac} % для "косых" дробей типа 1/2
\usepackage{makecell} % для ручного переноса строк в ячейках таблиц
\usepackage{datatool} % для работы с файлами данных